\section{Un prince au coeur de cœur de la foi chrétienne}

\subsection*{Introduction}

le simple fait d'être chrétien permet il de se trouver légitime face à l'Église, qustion qui intervient dans l'essence du pouvoir, sa provenance, mais aussi sa transmission à tarevrs la foi et qui baouti à un questionnement sur la façon dont Ambroise reágit face à un pouvoir qui ne défend pas la chrétienté. + questionner le fait pas seulement qu'il y ait le chrétienté mais que le prince soit chrétien et ce que ça implique.

\subsection{La seule chrétienté comme légitimité du pouvoir ?}

\subsubsection{L'institution du pouvoir par Dieu}

La réflexion ambrosienne sur le pouvoir impérial ne s'arrête évidemment pas à des comparatifs bibliques ou à une revalorisation chrétienne d'une pensée philosophique ancrée dans les mœurs romaines. Pour Ambroise, questionner le pouvoir et la figure de l'empereur, c'est donner la possibilité de le faire évoluer dans une direction voulue par l'évêque. Cet exercice de réflexion se fait notamment à travers la foi du monarque et à travers l'ensemble de la relation entre le pouvoir et la chrétienté. Ainsi, il est nécessaire d'analyser la façon dont l'évêque de Milan s'interroge sur la nature de l'autorité politique et du pouvoir. En effet, sa théorie politique d'une obéissance totale aux détenteurs du pouvoir et son rejet de toute forme de contestation physique, comme nous avons pu le voir précédemment, ne proviennent pas de rien, au contraire, ils s'inscrivent dans la continuité de la rhétorique de l'obéissance des premiers textes et penseurs chrétiens.

\bigskip

La raison de cette approche du pouvoir dans l'œuvre d'Ambroise s'explique par le fait qu'il s'appuie sur les fondements de la théologie chrétienne, avec l'Épître aux Romains de Paul : «~Que toute personne soit soumise aux autorités supérieures, car il n'y a point de pouvoir qui ne vienne de Dieu, et les pouvoirs qui existent ont été institués par Dieu.\footnote{Romains 13, 1.}~» Le cœur de cette réflexion est évidemment d'éviter toutes formes de révoltes des communautés chrétiennes pour s'assurer un chemin vers le Salut, malgré un contexte de persécutions de plus en plus important tout au long des trois premiers siècles de notre ère. L'enjeu n'était alors pas la justice temporelle de l'Empire, mais bien la justice de Dieu après la mort. Ambroise évolue dans un contexte politique bien différent avec un pouvoir désormais pleinement ancré dans la foi chrétienne, il peut continuer à porter la vision d'une «~donation~» du pouvoir par Dieu sans subir les contraintes d'un souverain persécuteur. Il réutilise donc fréquemment ce postulat de l'origine divine du pouvoir, mais insiste également sur l'aspect contraignant que cela implique, notamment la nécessité de se soumettre à la foi et à la piété pour conserver toute sa légitimité. Il peut entre autres s'appuyer sur l'Ancien Testament, dont nous avons déjà démontré l'importance dans l'ensemble de sa pensée, en rappelant l'histoire de la royauté d'Israël et l'autorité de Dieu dans le choix du pouvoir : Il punit Saül pour ses fautes et nomme David à sa place. Seule l'onction reçue par Dieu permet de s'emparer du pouvoir, ce qui amène également les souverains à être en partie tributaires de l'approbation divine. Plusieurs de ces aspects seront abordés dans la suite du chapitre, puisqu'il se concentre sur l'ensemble des notions chrétiennes touchant le pouvoir romain et la pensée d'Ambroise.

\bigskip

L'onction de Dieu est un principe très présent dans la pensée d'Ambroise, il s'adresse souvent aux empereurs en rappelant ce lien d'autorité qui les relie au divin :

\bigskip

\begin{quote}
«~[...], Seigneur empereur Auguste, élu par le jugement divin, très glorieux parmi les princes.\footnote{« \latin{[...], domine imperator auguste, divino electe iudicio, principum gloriosissime}.~». \cite[\nopp ec. 12, 10]{ambroise_csel_82_3}. Traduction personnelle.}~»
\end{quote}

\bigskip

Il y a par exemple, dans cette introduction de sa première lettre à Gratien, la précision importante de l'élection divine qui représente exactement l'onction, bien que ce soit un terme plutôt utilisé par les Écritures que par les penseurs chrétiens. Cette onction amène justement parfois une confiance abusive du monarque en son pouvoir : comme le souligne une étude sur l'héritage politique de la Bible, «~les souverains ont souvent abusé de leur lien particulier avec Dieu pour revendiquer un pouvoir absolu et une inviolabilité totale\autocite{mordecai_racines}~», alors qu'au contraire, l'idée même de l'onction est là pour rappeler l'autorité de Dieu sur le temporel. D'une certaine manière, en insistant sur le choix divin à l'origine du souverain, Ambroise affirme la légitimité de la présence du système monarchique dans le monde romain.

\bigskip

C'est dans son \latin{Expositio evangelii secundum Lucan} qu'on retrouve le plus d'informations sur la vision ambrosienne de l'origine divine du pouvoir. Ce commentaire de l'Évangile de Luc, certainement rédigé vers 389-390, revient sur la relation entre Dieu et l'exercice concret du pouvoir. Dans le livre IV, se pose un problème important de la pensée chrétienne : comment un pouvoir institué par Dieu peut-il être amené à s'opposer à la foi chrétienne, voire parfois à laisser apparaître un pouvoir tyrannique ? De nouveau, l'évêque de Milan reprend les fondements théologiques chrétiens, en séparant l'origine du pouvoir et son utilisation :

\bigskip

\begin{quote}
«~De Dieu vient l’institution des pouvoirs, du mauvais l’ambition du pouvoir.\footnote{« \latin{A deo potestatum ordinatio, a malo ambitio potestatis}.~». \cite[\nopp IV, 29]{ambroise_csel_32_4}.}~»
\end{quote}

\bigskip

Par cette formule, Ambroise n'inclut pas Dieu dans les fautes impériales, ce qui permet de donner une explication au fait que Dieu, dans l'histoire du christianisme, ait pu laisser le pouvoir entre les mains de persécuteurs. Il ne s'agit pas d'un mauvais choix, mais d'une corruption de l'humain par le diable, ce qui est facilité par la possession du pouvoir et de l'autorité politique. On peut d'ailleurs noter que cette précision rejoint la critique, évoquée à plusieurs reprises dans le premier chapitre, envers l'ambition personnelle qui représente selon lui l'aspect le plus propice à aboutir à la corruption de la société.

\bigskip

Cependant, bien que les fautes du souverain ne doivent pas être attribuées à Dieu, Ambroise se permet d'établir un lien entre les actions du monarque envers la foi et son statut au cœur de la religion. C'est ce que l'on retrouve dans la suite immédiate de son commentaire de Luc :

\bigskip

\begin{quote}
«~Au reste, l’institution du pouvoir vient tellement de Dieu que celui qui use bien du pouvoir est ministre de Dieu. [...] Ce n’est donc pas la fonction qui peut être en faute mais le ministre, ce n’est pas l’institution de Dieu qui peut déplaire, mais la conduite de celui qui administre.\footnote{« \latin{Ideo potestas a deo est, ut is qui bene utitur potestate minister dei sit. [...] non ergo munus in culpa est, sed minister, nec dei potest ordinatio displicere, sed usurpantis administratio}.~». \cite[\nopp IV, 29]{ambroise_csel_32_4}.}~»
\end{quote}

\bigskip

L'utilisation du terme « ministre de Dieu » est surprenante car elle relie la figure du souverain, de l'empereur, aux évêques qui sont les véritables ministres de Dieu. Bien que tout pouvoir vienne de Dieu, l'exercice de ce pouvoir diffère grandement d'un monarque à un autre. C'est donc par leur implication dans les affaires religieuses et par le suivi du chemin voulu par Dieu que se démarquent certains souverains, qui peuvent espérer une place plus importante dans la vie religieuse. Mener une politique terrestre qui correspond aux exigences divines n'est pas un prérequis du pouvoir, mais permet au souverain de s'assurer son propre Salut tout en garantissant une certaine stabilité à l'ensemble de la société, et donc de l'Empire.

\bigskip

Ce fondement théologique permet d'expliquer la cohérence de l'ensemble de la réflexion ambrosienne. C'est parce que l'institution du pouvoir est offerte par Dieu qu'il peut prôner le respect du pouvoir en place et s'opposer avec virulence aux guerres civiles et aux usurpateurs. Et on comprend avec la dernière citation que l'onction divine est ce qui justifie sa valorisation des empereurs qui luttent contre le polythéisme romain et les courants hérétiques du christianisme. Le pouvoir en lui-même est légitime peu importe l'utilisation qui en est faite, mais ne peut être considéré comme «~bon~» par Ambroise que si l'empereur est ministre de Dieu.

\bigskip

Enfin, cette lecture importante sur la façon dont Ambroise parle de l'origine du pouvoir nous amène à revenir sur ce qui peut s'apparenter à une contradiction : son questionnement sur le système hiérarchique et, plus généralement, sur la façon de nommer le futur souverain\footnote{Voir le 2.1.3 du Chapitre 1.}. Comme nous l'avons analysé, l'évêque de Milan rejette la désignation par le hasard, préférant la recherche d'un monarque qui se détache par les attraits de la nature, à l'image de la reine dans la ruche, de taille bien plus imposante. Ce choix par la distinction naturelle se révèle finalement être l'expression de ce don du pouvoir par Dieu, l'élection divine ne faisant que confirmer l'ordre naturel déjà présent. Cela n'empêche évidemment pas la présence d'une succession dynastique, tant qu'elle reflète le choix de Dieu, ce qui n'est pas le cas des usurpateurs, même acclamés par l'armée et se revendiquant de foi chrétienne.

\bigskip

\subsubsection{Légitimer la descendance par la foi du souverain}

L'origine divine du pouvoir justifie donc l'exercice du pouvoir impérial dans sa globalité, tout en valorisant les souverains qui agissent en faveur de la foi chrétienne. Mais cet aspect n'est pas le seul à être pris en compte par Ambroise pour consolider sa construction d'un empereur chrétien idéal. La réflexion de celui-ci s'attarde sur l'ensemble des implications, positives comme négatives, autour du fait que le détenteur du pouvoir soit chrétien. Et intervient alors notamment la question de la place de la foi du souverain dans la succession et la descendance. Comme nous l'avons précédemment analysé, l'idée d'un système purement dynastique est théoriquement critiquée par Ambroise, qui y voit un risque important pour la stabilité de la société. Pourtant, dans la pratique, l'évêque de Milan accepte, voire soutient, le principe de succession familiale lorsqu'intervient un principe fondamental : une foi suffisante lors de la vie du père défunt.

\bigskip

Son acceptation du principe dynastique est liée à un élément important que l'on retrouve dans cette succession : la piété du père. Ambroise ne devient pas un défenseur de l’hérédité du pouvoir, mais il l'accepte dans le cadre d’une succession suivant les valeurs de fidélité à Dieu, la fameuse notion de ministre de Dieu que nous venons d'évoquer. Il reste ainsi cohérent avec ses fondements sur l'origine divine du pouvoir et avec la possibilité d'entretenir un lien avec Dieu à travers la foi de l'empereur. Ce mécanisme de légitimation spirituelle se retrouve tout particulièrement dans la succession de Théodose avec ses deux fils Arcadius et Honorius. Tout au long de son oraison funèbre, le \latin{De obitu Theodosii}, Ambroise prend le temps d'associer les deux nouveaux Augustes à leur père, en rappelant les vertus et la foi de ce dernier pour ancrer leur nouvelle autorité politique dans une continuité morale et spirituelle et leur garantir une forme de réussite terrestre.

\bigskip

Le lien de l'empereur avec le divin n'est pas nouveau, mais c'est la forme qui désormais diffère. Dans la tradition romaine, l'empereur défunt est divinisé et apporte une protection à sa descendance par ce statut de \latin{divus}. L'arrivée du christianisme chez les empereurs bouleverse cette approche puisqu'il n'est plus envisageable de faire un rapprochement entre le monarque et Dieu. La réussite de la descendance ne dépend alors plus du statut divin du père, mais se concentre désormais sur les actions de foi, la fameuse piété du prédécesseur. L'accumulation des actions en faveur de la foi chrétienne permet d'assurer une continuité dans les règnes et donc facilite la stabilité de la société romaine :

\bigskip

\begin{quote}
«~Théodose au contraire était rempli de la crainte de Dieu, plein de miséricorde. Nous espérons qu’il est, auprès du Christ, un protecteur pour ses enfants, si le Seigneur est favorable aux affaires humaines.\footnote{« \latin{Theodosius uero plenus timoris Dei, plenus misericordiae. Speramus quod liberis suis apud Christum praesul adistat, si Dominus propitius sit rebus humanis}.~». \cite[\nopp 16]{ambroise_mort_theodose}.}~»
\end{quote}

\bigskip

La continuité politique que recherche Ambroise dans le contexte de succession de l'Empire trouve sa réussite dans la protection que l'empereur pieux peut apporter à ses enfants, ici dans l'exemple de Théodose. Pour forger la figure chrétienne du souverain, l'évêque joue avec les proximités familiales pour amener la conscience des nouveaux empereurs à suivre les prédécesseurs valorisés par Ambroise. C'est donc le cas avec les fils de Théodose, mais c'est ce que l'on retrouve également dans les lettres à destination de Valentinien II : il montre au jeune empereur qu'il ne peut pas prendre des mesures allant contre la chrétienté puisque son père et son frère, Valentinien et Gratien, étaient pleinement ancrés dans la défense de la foi et faisaient preuve d'une grande piété\autocite[\nopp 72, 16-17]{ambroise_lettres}. On a dans cette citation à la fois l’idée d'intercession de l’empereur à destination de ses héritiers, mais également la présence de cette intervention divine sur Terre sous différentes formes. Un peu plus tôt dans le discours, l’évêque utilise d’ailleurs le terme «~Grâce au Seigneur~», continuant dans cette idée d’une aide de Dieu en cas de bonne conduite dans la foi, aide qui légitime la succession dynastique. Nous reviendrons plus en détail sur le concept de Providence divine dans la suite du chapitre.

\bigskip

Le facteur le plus important que veut mettre en avant Ambroise, et qui explique en grande partie son soutien à la prise du pouvoir d’Honorius et Arcadius, est l’idée que la piété du prince régnant influe directement sur sa descendance et donc sur l’héritage politique qu’il laisse. Pour Ambroise, le fait que Théodose soit resté dans la foi tout au long de son règne est signe du futur succès des jeunes Auguste. Cette pensée mettant en évidence une relation entre la piété et la réussite d’une même famille est évidemment de nouveau issue des textes bibliques, et principalement des livres de Samuel et des Rois, fondement de la théorie politique d’Ambroise. Il défend déjà cette idée dans son \latin{De Officiis} en parlant de la grande piété de David et de tous ses accomplissements pour la foi qui ont influencé la suite de la politique du royaume d'Israël :

\bigskip

\begin{quote}
«~Lui à cause de qui le pardon fut accordé à ses descendants pour leurs fautes et à cause de qui son privilège fut conservé à ses héritiers.\footnote{« \latin{texte en latin}.~». \cite[\nopp II, 35]{ambroise_devoirs_2}.}~»
\end{quote}

\bigskip

La vertu du défunt peut donc apporter un soutien providentiel direct à sa lignée, allant jusqu'à effacer les propres fautes de la possible descendance. Il est ici question de Salomon, dont Dieu conserve l'unité du royaume malgré les fautes personnelles du jeune roi, uniquement en l’honneur des actions passées de David. La comparaison prend tout son sens : comme pour Salomon, la légitimité d’Honorius et Arcadius provient du fait qu'ils succèdent à un empereur pieux et soutenu par Dieu. Le développement de cette rhétorique ne se fait pas au hasard, Ambroise cherche à valoriser la piété des empereurs en la reliant à un concept très romain de la filiation et la continuité familiale. L'évêque est conscient en 395 de la faiblesse potentielle que représente le jeune âge des deux héritiers, il est donc nécessaire de rassurer les élites et d'avorter les possibles contestations en donnant une légitimité religieuse à leur présence. Le modelage du prince chrétien continue de se faire entre idéal et pragmatisme, et la question de la foi pour la descendance relie les deux. Défendre la légitimité d'Honorius et Arcadius sur le terrain de la foi a donc le double objectif de renforcer le pouvoir des jeunes empereurs ainsi que les obliger à suivre plus fidèlement encore la foi chrétienne et les conseils du monde épiscopal.

\bigskip

Pour appuyer son propos, Ambroise utilise dans son oraison funèbre plusieurs figures bibliques qui confirment l’importance de la foi dans la transmission du pouvoir. Il se permet ainsi d'utiliser son sermon comme un enseignement complémentaire pour l’ensemble de son auditoire. Après avoir décrit le règne de Josias, aussi bien sa piété que ses erreurs, il développe le cas de Asa, jeune roi d'Israël qui parvient à rester au pouvoir pendant quarante ans. Fidèle à Dieu et dans une logique de piété, il demande de l’aide divine alors qu’il se bat contre des Éthiopiens bien plus nombreux, ce qu’il obtient. Mais, lorsqu’il tombe malade, Asa « délaisse » la foi et la confiance en Dieu en demandant des soins à des médecins syriens n’étant pas dans la même religion. La tradition biblique fait de cette erreur le facteur expliquant sa mort dans la souffrance. Mais c’est surtout la deuxième partie de l’exemple d’Ambroise qui lui permet de justifier son idée de l’importance de la foi du père pour le fils. En effet, Ambroise creuse l’exemple de Josias et Asa et justifie les erreurs des deux rois par l’absence de foi de ceux dirigeant avant eux :

\bigskip

\begin{quote}
«~Mais leurs pères, Abiyam et Amos, étaient tous deux des hommes sans foi en Dieu.\footnote{« \latin{Sed illorum patres, Abias et Amos, infideles}.~». \cite[\nopp 16]{ambroise_mort_theodose}.}~»
\end{quote}

\bigskip

En plus de la foi des enfants, bien évidemment nécessaire à la réussite du règne, Ambroise fait intervenir la piété des parents comme facteur décisif dans la prospérité du royaume et de la vie terrestre des princes. Le message que cherche à faire passer Ambroise dans l’ensemble de son oraison, et plus généralement tout au long de sa proximité avec le pouvoir, est donc clair : si la foi du père protège la dynastie, comme l’exemple de David et ses descendants, la négligence dans la piété la condamne au malheur du règne terrestre avec l’absence du soutien divin.

\bigskip

Après avoir longuement loué les vertus religieuses de Théodose, Ambroise place donc définitivement sa conduite comme un gage de continuité de la protection divine, de sécurité et de prospérité à venir pour ses deux enfants. Ainsi, dans le contexte précis d’Honorius et Arcadius, Ambroise ne rejette pas le principe dynastique mais cherche à l’inscrire dans une logique de continuité spirituelle entre leur foi et celle de Théodose. Il continue de s’opposer à la seule légitimité de la naissance mais sait adapter son discours à la situation qu’il doit résoudre pour assurer la stabilité de l’Empire et confirmer son propre rôle auprès du pouvoir, deux logiques fondamentales chez Ambroise et que l’on retrouve clairement dans son oraison funèbre à Théodose. Ces analyses nous permettent de comprendre avec précision le regard d'Ambroise sur l'aspect chrétien du prince, pour réussir à pousser le pouvoir vers la réalisation de son idéal d'intérêt général et de morale chrétienne.

\bigskip

\subsubsection{Les limites du loyalisme politique}

La foi chrétienne des détenteurs du pouvoir semble offrir un socle de légitimité suffisant pour l'Église et Ambroise, mais cet aspect peut également être sujet à des critiques lorsque le pouvoir s'écarte de la voie tracée par la chrétienté qu'Ambroise cherche constamment à consolider. À ce sujet, la pensée politique de l'évêque de Milan repose entre autres sur une tension constante entre la fidélité à l'autorité impériale, justifiée par son origine divine, et la primauté absolue de la foi chrétienne. Toute la réflexion ambrosienne autour de la légitimité et du loyalisme se tourne vers ce questionnement : si l'empereur cesse de défendre la foi chrétienne, quelle doit être l'attitude de l'Église et celle d'Ambroise ? Le respect du pouvoir porté par l'évêque n'est pas non plus le signe d'une soumission aveugle à toute forme d'autorité terrestre. De nouveau il peut conforter sa position à partir des préceptes bibliques : il convient avant tout d'obéir à Dieu plutôt qu'aux hommes\footnote{Actes 5, 29.}, ce qui est pris par Ambroise comme le symbole du fait de s'écarter des actions politiques allant contre Dieu et contre le christianisme. Cette limite théorique trouve son application dans la vie d'Ambroise lors de l'usurpation d'Eugène en 392, qui pose à l'évêque de Milan nettement plus de problèmes politiques et théologiques que l'usurpation de Maxime et la guerre civile qui s'en est suivie.

\bigskip

La chronologie du conflit opposant Théodose et Ambroise à Eugène est également importante à prendre en compte pour saisir pleinement les actions et la pensée d’Ambroise. À la mort de Valentinien II, qui décède sans héritier, un vide politique apparaît en Occident, car le seul empereur légitime encore présent, Théodose, souhaite rester dans la partie orientale de l'Empire, ce qui amène des contestations quant à sa capacité à régner sur un territoire éloigné\autocite[343]{mclynn_ambrose}. Celui-ci souhaite logiquement imposer son fils Honorius en tant qu'Auguste, mais son très jeune âge ne lui apporte pas une légitimité suffisante pour s'imposer en Occident. Eugène, soutenu par le général franc Arbogast, a alors le champ libre pour s'imposer sans faire face à trop de contestations puisqu'un pouvoir fort semble nécessaire pour garantir la sécurité des frontières de l'Empire. Théodose s'empare tout de même du contrôle de l'Italie, mais laisse les provinces les plus à l'ouest à Eugène. Face à cette situation, la position d'Ambroise n'apparaît pas clairement, puisqu'il souhaite suivre le silence diplomatique de Théodose. En effet, l'empereur reste vague : il reçoit les ambassades d’Eugène mais sans donner suite et décide de ne pas légitimer l’usurpateur, sans pour autant l’affronter tant que celui-ci reste en Gaule, territoire qu’il ne convoite pas, y compris pour ses fils\autocite[266-270]{maraval_theodose_2009}. N'ayant plus de soutien fort en Occident avec la mort de Valentinien, Ambroise préfère ne pas rentrer dans un conflit ouvert et s’aligne sur la position de Théodose, il s’efface des affaires politiques et décide de ne pas répondre aux différentes demandes d’Eugène.

\bigskip

Mais toute la réflexion sur le loyalisme politique d'Ambroise, et sur sa volonté de guider les empereurs vers un idéal de la monarchie chrétienne, s'accélère lorsque Théodose nomme officiellement Honorius en Italie : Eugène brise le flou politique en rentrant en conflit ouvert avec Théodose en amenant son armée sur le sol italien. Face à l'échec de ses ambassades auprès de Théodose, il est conscient qu'il ne sera pas reconnu officiellement. Il cherche alors du soutien parmi l'aristocratie romaine, et trouve notamment du soutien dans le monde polythéiste. Bien qu'étant lui-même chrétien, il se lance alors dans la promesse d'une restitution des cultes païens pour s'assurer la présence de soutiens puissants dans sa conquête de l'Italie et du pouvoir. Il ne mène pas une véritable politique anti-chrétienne, mais prône une tolérance qui déplaît aux chrétiens d'Italie et notamment à Ambroise. Cette politique couplée à sa volonté de prise de l'Italie permet à Ambroise de se libérer dans la position qu'il souhaite adopter, et montre alors une véritable limite à son loyalisme politique : face à un pouvoir agissant contre la chrétienté, l'évêque refuse de s'y soumettre. Il refuse de reconnaître l'autorité politique d'Eugène, mais refuse également de le confronter. Il quitte ainsi Milan avant d'en informer Eugène pour lui expliquer sa position :

\bigskip

\begin{quote}
«~Je me suis retiré de Milan par crainte de Dieu, vers qui j’ai l’habitude de diriger tous mes actes autant que je le peux ; jamais je ne détourne mon esprit de lui, et je ne tiens pour plus précieuse la faveur d’aucun homme que celle du Christ.\footnote{« \latin{Secessionis mihi causa timor domini, ad quem omnes actus meos quantum queo dirigere consuevi nec umqua ab eo mentem deflectere nec pluris facere cuisvis hominis quam Christi gratiam}.~». \cite[\nopp ec. 10, 1]{ambroise_csel_82_3}.}~»
\end{quote}

\bigskip

Concernant ce retrait volontaire, McLynn propose une nouvelle analyse. Pour l'historien, la justification de son départ par crainte de Dieu n'est qu'un élément de rhétorique qui est en réalité un mensonge servant à conforter son image au sein de la communauté chrétienne\autocite[345-346]{mclynn_ambrose}. McLynn estime qu'il utilise cet argument pour justifier un départ qui en réalité lui permet simplement de ne pas se mettre à dos Théodose, tout en évitant un véritable conflit avec Eugène. L'évêque profite d'ailleurs de cette absence pour engager un cheminement politique et religieux en Italie du Nord, renforçant son pouvoir épiscopal plutôt que luttant contre la politique religieuse d'Eugène. Bien que le ton de la lettre, loin d'être agressif, aille dans le sens de l'analyse de McLynn, il ne faut tout de même pas nier l'engagement de l'évêque dans la lutte contre le polythéisme ni oublier ses prises de positions contre Valentinien ou Théodose lors des actes politiques n'allant pas dans le sens de la vision de la politique religieuse d'Ambroise. Il est donc évident que son pragmatisme l'amène à faire attention à sa position dans la politique impériale, mais il existe tout de même une opposition théologique à Eugène.

\bigskip

Ce geste évoqué dans la citation précédente symbolise bien la pensée d'Ambroise : Il ne montre pas un acte de rébellion mais plutôt une forme de résistance passive envers une autorité ne respectant pas la foi chrétienne. Il se légitime ainsi par la primauté de l’autorité de Dieu, rappelant le lien permanent qu’il fait entre légitimité du pouvoir et foi en Dieu. Si l’autorité divine est remise en cause par les actions du monarque, l’Église a le devoir de s’en écarter. Ambroise se permet même de conclure sa lettre à Eugène de façon assez agressive sur ce sujet de hiérarchie entre empereur et Dieu :

\bigskip

\begin{quote}
«~Mais puisque vous désirez que l'on vous rende hommage, permettez-nous de rendre hommage à celui de qui vous souhaitez prouver que votre autorité provient.\footnote{« \latin{Sed qui vobis deferri vultis patimini, ut deferamus ei quem imperii vestri vultis auctorem probari}.~». \cite[\nopp ec. 10, 12]{ambroise_csel_82_3}.}~»
\end{quote}

\bigskip

L'évêque de Milan s’oppose ainsi frontalement à Eugène ainsi qu’à l’ensemble de la politique qu’il cherche à mener en Occident. La lettre d’Ambroise ne signifie pas pour autant une rupture définitive avec le pouvoir impérial, il continue d’utiliser les formes de courtoisie et de respect notamment par le « très clément empereur » que l’on voit à plusieurs reprises. En revanche, il n’y a pas dans sa rédaction de volonté d’aide spirituelle. Il n’essaie pas, comme il a pu le faire avec Théodose après les événements de Callinicum ou de Thessalonique, de « sauver » Eugène en le forçant à réparer ses fautes. Le symbole de cette différence entre les correspondances avec les détenteurs du pouvoir se voit notamment par l’absence des références aux rois bibliques dans la lettre à Eugène, pourtant si fréquemment utilisées par Ambroise. L’usurpateur n’apparaît donc pas aux yeux d’Ambroise comme un empereur fautif, mais bien comme un monarque sans légitimité spirituelle, ce qui est confirmé par sa lettre à Théodose après sa victoire à la bataille de la Rivière Froide où l’évêque évoque Eugène en le qualifiant directement d’«~usurpator indignus~», un usurpateur indigne\autocite[\nopp ec. 2, 1]{ambroise_csel_82_3}.

\bigskip

L’épisode de l’épiscopat d’Ambroise avec Eugène reste donc au cœur de sa ligne politique : il reste fidèle à son modèle de loyalisme et de respect de l’autorité, qu’il met régulièrement en avant par le biais de David avec Saül, mais rappelle que l’obéissance n’est un devoir que lorsque le pouvoir se soumet à la volonté divine, ce qui n’est pas le cas d’Eugène qui s’allie avec les figures païennes et favorise le retour de la religion romaine traditionnelle. Il ne cherche pas pour autant la violence qui engendre la guerre civile, mais il s'écarte de ce nouveau pouvoir et met en place une critique morale importante. Cette analyse permet une continuité avec l'importance de la chrétienté comme origine du pouvoir et comme facteur de légitimation de la descendance. Pour Ambroise, le seul modèle monarchique qui peut aboutir à la réussite de la société est lorsque le pouvoir est conforme à la morale chrétienne et suit les conseils de l'idéal ambrosien.

\bigskip
